\section{Conclusion}
\label{sec:conclusion}

This paper introduces a fully deployable intrusion detection system for satellite Controller Area Network (CAN) bus networks, requiring only 541 bytes of flash memory and executing in 4.2 microseconds on a STM32F373C8T microcontroller at 72 MHz. The system leverages 14 domain-specific CAN features and implements a depth-5 decision tree as static arrays, thereby eliminating all runtime dynamic allocation and ensuring deterministic, bounded inference, as validated across 348,000 hardware-classified frames. Protecting satellite CAN networks is directly relevant to energy infrastructure security: satellites providing GPS timing synchronization, renewable energy telemetry relay, and grid monitoring communications represent a single point of failure whose compromise can cascade across multiple energy sectors, as demonstrated by the 2022 Viasat attack on Enercon wind turbine operations~\cite{Viasat2022, SentinelOne2022}.

Live SATLL Electrical Power System (EPS) telemetry quantifies the intrusion detection system (IDS) power consumption at two levels. The inference increment, defined as the marginal algorithmic cost at 100 Hz, is bounded at \textbf{48--53~\textmu W} ($\Delta P \leq I_\text{base} \cdot V_{dd} \cdot t_\text{inf} \cdot f_\text{poll}$, measured $I_\text{base}=35$--$38$~mA). This represents less than 0.004\% of the 1,273–1,471 mW added by the Attitude Determination and Control System (ADCS) subsystem during sensor experiments. The current evaluation board draws 164--176~mW from the satellite payload port (MCU base 115--125~mW, USB bridge and peripherals $\approx$49~mW); this overhead belongs entirely to the evaluation hardware, not the algorithm. The external board also adds mass, a practical constraint for CubeSat-class platforms where every gram affects launch cost and orbital lifetime. The most impactful path to production deployment is direct integration of the IDS into the STM32H573 OBDH firmware, which would eliminate the external board entirely and reduce the marginal IDS cost to $\approx$40~\textmu W. Direct integration was not attempted in this work because the SATLL OBDH runs closed-source firmware that does not permit third-party modification. Our results demonstrate that the algorithm itself imposes negligible overhead; future work will pursue OBDH-integrated deployment on an open-firmware platform.

The cross-domain evaluation quantifies the generalization gap observed when the satellite-trained model is applied to unfamiliar network traffic, underscoring the need to collect authentic mission traffic for production training. The current training corpus covers four canonical attack classes; operational deployment will encounter additional traffic patterns---eclipse transitions, orbit corrections, payload activations---not represented here, which may increase false positives. A static depth-5 tree also presents a limited decision surface: an adversary with feature-set knowledge could craft evasion frames, motivating future hardening within the same flash budget.

\section*{Reproducibility}
All training code, feature engineering scripts, model export toolchains, and benchmark scripts are included with the paper submission. Trained model files and C header artifacts are also provided to facilitate direct replication of the firmware validation results.